\section{Introduction}
\label{sec:introduction}

Every claim a modern autonomous-vehicle perception stack makes about the world -- that box is a pedestrian, that lane is clear, that gap is wide enough to merge into -- is implicitly conditioned on the sensing conditions the detector was trained and validated under. Snow breaks that conditioning quietly. Unlike a sensor failure or a dropped frame, snowfall does not announce itself to the software stack: airborne flakes register on LiDAR as spurious near-range returns rather than as an error code, camera contrast degrades gradually rather than cutting to black, and a 3D detector that was accurate a minute ago keeps producing confident, plausible-looking boxes as its actual reliability erodes underneath it. The failure mode that matters operationally is not "the detector is wrong" -- that is always somewhat true -- but "the detector has become wrong in a way the system has not been told about." Detecting that transition, with a false-alarm rate the vehicle's safety case can actually budget for, is a different problem from building a more snow-robust detector, and it is the problem this paper addresses.

Conformal prediction offers a distribution-free way to attach calibrated uncertainty to a black-box detector's outputs: calibrate a nonconformity score on held-out data, and a new prediction's miscoverage is controlled at a chosen level regardless of what the detector is or how it was trained. That guarantee, however, is a batch guarantee. Applied naively to a live video stream -- recomputing a conformal test at every frame, or every fixed window of frames -- it either inflates the true false-alarm rate under continuous monitoring (the classic multiple-testing failure of checking a hypothesis test more than once) or requires a Bonferroni-style correction so conservative that real onset events are missed by the time the correction lets the alarm fire. What a deployed vehicle actually needs is a test that stays valid \emph{no matter when or how often it is checked} -- anytime-valid, in the sequential-testing sense -- so that "has anything changed yet?" can be asked continuously without paying a multiple-testing tax and without fixing a monitoring horizon in advance.

Sequential testing by betting supplies exactly that guarantee. Frame the question as a bet: under the null hypothesis that the detector is still operating within its calibrated miscoverage budget, a wealth process built from the observed miscoverage rate is a nonnegative supermartingale, and Ville's inequality bounds the probability that this wealth ever crosses an alarm threshold, at any stopping time, not just at a horizon fixed in advance. This backbone -- e-processes, testing by betting, Ville's inequality -- is well established in the sequential-testing literature~\cite{Ramdas2023GameTheoretic,WaudbySmithRamdas2024JRSSB}, and it is not, on its own, this paper's contribution. The direct predecessor for applying it to a live perception problem is Monroy Mu\~noz, Verma, and Timans's WACV~2026 work on sequential object-tracking-failure detection~\cite{MonroyMunoz2026WACV}: an e-process over a bounded tracking-quality metric, with the betting rate learned adaptively from the metric's own history via aGRAPA or SF-OGD, evaluated across four 2D visual-tracking benchmarks. It is a strong, directly applicable baseline, and we build on it rather than around it -- everywhere in this paper that we compare against a "covariate-blind" betting rule, we mean their rule, not a strawman.

What their formulation leaves open, and what we take as our starting point, is twofold, plus a third direction we explored and report honestly rather than claim. First, their evaluation is explicitly single-object and per-video; a driving scene contains many objects and, more importantly, spatial structure that a single global e-process discards entirely -- the perception system can be untrustworthy in the near lane and fine on the shoulder, and a monitor that only ever says "something, somewhere is wrong" is far less actionable for a downstream planner than one that says where. Second, their benchmarks are all 2D visual tracking under whatever distribution shift those datasets happen to contain; ours is 3D multimodal driving perception under a real, physically grounded weather transition on a dataset with no clear-weather baseline at all, which is a different failure surface with different evaluation challenges than theirs. Third, their betting rate is learned purely from the failure metric's own past values; nothing in the update conditions on an independent signal that might move \emph{before} the failure metric does, and a multimodal 3D detector exposes exactly such a signal for free if it fuses camera and LiDAR through some cross-modal agreement mechanism -- we tried conditioning on it, and Section~\ref{sec:experiments} reports, in full, why that attempt did not work.

Concretely, our primary contributions target the first two gaps. We replace the single global wealth process with a grid of per-BEV-cell wealth processes, one per spatial region, controlled jointly under an online false-discovery-rate procedure (e-BH~\cite{WangRamdas2022JRSSB}) so that testing many regions simultaneously does not inflate the overall false-alarm budget, turning one alarm bit into a live map of where the stack is currently untrustworthy -- Section~\ref{sec:method} gives the exact construction. And we build and report the first evaluation of this monitoring approach, to our knowledge, against a real adverse-weather driving dataset rather than only synthetic corruption of a clear-weather one: Snowy Scenes, captured entirely in genuine snowy conditions, which turned out to contain no clear-weather split at all, forcing a real severity-ordering methodology and a real (if constructed) onset-stream splice in its place -- Section~\ref{sec:experiments} explains the dataset, the construction, and the resulting real operating curve, on which the covariate-blind backbone itself detects the real weather onset within 3 to 10 frames depending on the false-alarm budget, with zero false alarms across every budget we tested.

Alongside those two contributions, we also tried conditioning the betting rate on the disagreement signal from the detector's own Cross-modal Consistency Probe (CCP), so the process could bet more aggressively on frames where the two modalities already disagree rather than waiting for the lagging miscoverage rate to catch up -- an appealing idea that, on real data, ran into two distinct real problems rather than one. First, the trained CCP score collapsed to near-constant agreement, because nothing in the detector's own training objective had ever taught it to disagree; we diagnosed and fixed that training-objective gap. Second, once fixed, the score varied again but in the wrong direction for this task, and the covariate-informed bettor's real operating curve came out worse, not better, than the covariate-blind baseline's. We report both findings in full in Section~\ref{sec:experiments}, because we think a well-optimized internal signal becoming a misleading external covariate -- not once, but in two different ways across two different fixes -- is a useful cautionary result for the next paper that tries this, and because burying a negative result would misrepresent what this paper's evidence actually supports. We treat covariate-informed betting via this specific mechanism as an open problem this paper investigated and could not close, not as a third headline contribution.

We are correspondingly careful about what we claim for the two contributions we do make. The evaluation in this paper is a first real demonstration at a modest scale -- a handful of replicate onset streams, short monitoring windows, one detector architecture, one dataset -- not a large benchmark, and Section~\ref{sec:conclusion} is explicit that the spatial multiplicity-control grid, while implemented and validated end to end on the real detector and dataset, has not yet been run against a real multi-object driving scene with genuine spatial heterogeneity, which is the setting it is actually motivated by. Our contribution is the spatial monitoring construction, the real-data evaluation protocol and dataset access story, and -- reported with the same rigor as the rest of the paper rather than as an afterthought -- two documented, mechanistically diagnosed failure modes of the covariate-informed betting idea we also investigated. We do not claim the backbone e-process itself is novel, which it is not, and we do not claim covariate-informed betting works, which our own evidence argues against.
